% framing_task_regime_generalization.tex — 2026-08-01
%
% Standalone candidate replacement for the "Task-Regime Generalization"
% subsection in results_total/paper/EMNLP_My__Copy_/latex/sections/
% 4_experiments.tex (around line 288-300). NOT wired into that build --
% saved here for the author to review/paste in.
%
% Why this rewrite: the existing paragraph there claims Macro&Micro
% "achieves the best AUROC in every held-out regime" and "transfers
% consistently across task regimes." Checked against the current
% (2026-08-01, post outer-CV-shuffle-fix) numbers in
% code/generalization/results/cross_category_{results,summary}.csv and
% that claim does not hold: Macro&Micro wins AUROC on only 2/5 held-out
% categories, essentially tied with HTC (M11) on the mean (0.789 vs
% 0.787). The real, consistent advantage is on e-AURC/Brier (selective
% ranking + calibration), not raw AUROC. Confirmed this reading from
% the regenerated figure (results/website_generalization_figure.png,
% also updated 2026-08-01: title/subtitle removed, 3 weakest baselines
% dropped, HTC/Stacked_L2/Cross-Category Ours - M&M relabeled).
%
% Source data: code/generalization/results/cross_category_results.csv,
% cross_category_summary.csv (method column: ours_l2_all,
% m11_htc_full_step; category column: devtools, ecommerce,
% health_services, social_entertain, travel_navigate).

\subsection{Task-Regime Generalization}
\label{subsec:task_regime_generalization}

\begin{figure*}[t]
    \centering
    \includegraphics[width=\linewidth]{fig/website_generalization_figure.pdf}
    \caption{Task-regime generalization under key-step labels. We hold out
    one website category for testing and train on the remaining four.
    Cross-Category Ours - M\&M (red) is compared against HTC (orange, M11),
    Stacked\_L2 (M13), and the muted single-signal baseline cloud, across
    AUROC, e-AURC, and Brier score.}
    \label{fig:task_regime_generalization}
\end{figure*}

Figure~\ref{fig:task_regime_generalization} shows a more qualified picture
than raw discriminative power alone would suggest. On \textbf{AUROC}, the
pooled Macro\&Micro predictor is essentially tied with HTC: mean AUROC
across the five held-out categories is \(0.789\) vs.\ HTC's \(0.787\), and
Macro\&Micro wins only \(2/5\) categories outright (DevTools \(0.807\) vs.\
\(0.759\); Health Services \(0.787\) vs.\ \(0.752\)), losing on E-Commerce
(\(0.764\) vs.\ \(0.815\)) and Travel/Navigate (\(0.802\) vs.\ \(0.812\)),
and nearly tying on Social/Entertainment (\(0.787\) vs.\ \(0.794\)).
Macro\&Micro's advantage is not "higher AUROC in every regime" -- that
claim does not hold in the current data.

Where the dual-level representation \emph{is} consistently ahead is
\textbf{e-AURC and Brier score} -- metrics that penalize how reliably risk
is ranked and calibrated, not just whether it is ranked correctly on
average. Macro\&Micro wins e-AURC on \(3/5\) categories (mean \(0.123\)
vs.\ HTC's \(0.135\)) and Brier on \(4/5\) (mean \(0.187\) vs.\ \(0.193\)),
with a tighter per-category range on both (e-AURC \(0.113\)--\(0.143\),
Brier \(0.180\)--\(0.199\)) than HTC's. We read this as evidence that the
dual-level macro+micro representation's advantage over a purely
generation-confidence-derived signal like HTC lies in
\emph{selective-ranking reliability and calibration under category shift},
not in raw separability -- a narrower and more defensible claim than
"transfers consistently across task regimes." One plausible (unverified)
mechanism: HTC's score is a fitted probability over token-level
generation-confidence statistics, which may separate successes from
failures about as well as Macro\&Micro on average, but is less directly
probability-like as a raw signal, which could show up as worse
calibration/selective-ranking behavior under distribution shift even when
raw discrimination is comparable. This has not been independently
verified and should be treated as a hypothesis, not a conclusion.
